\section{Introduction}
\label{sec:introduction}

Tool-using agents interleave model decisions with retrieval, API calls, code
execution, and interactions among agents. When an execution fails, the final
error rarely identifies the decisive cause. A malformed argument may be
repaired downstream, an early context error may surface only after several
valid actions, and a failed tool call may be a consequence rather than the
cause. Agent failure diagnosis must therefore recover a causal report from the
trajectory and ground that report in evidence that another process can locate
and check.

Recent systems make trajectory-based attribution increasingly concrete.
AgentRx derives and checks execution constraints before localizing a critical
step~\cite{barke2026agentrx}; AgenTracer combines controlled injection,
counterfactual replay, and learned attribution~\cite{zhang2025agentracer}; and
VerifyMAS verifies failure hypotheses over complete multi-agent
trajectories~\cite{qiao2026verifymas}. Who\&When Pro further scales controlled
failure injection and golden attribution labels~\cite{liu2026whowhenpro}.
These methods improve localization, supervision, and hypothesis testing, but
they leave a downstream decision unresolved: when should a system trust a
structured diagnosis rather than treat it as another plausible model output?

Three failure modes make this trust decision distinct from attribution alone.
First, citation validity is weaker than evidential support: a model can cite a
real field that is irrelevant, insufficient, or contradicted elsewhere in the
trajectory. Second, a critic does not create independent evidence. A diagnoser
and verifier can share a model family, provider, prompt assumptions, retrieved
summaries, or the same misleading trajectory abstraction. Third, deployment
requires selective acceptance. Rejecting uncertain reports can lower observed
error only by sacrificing acceptance, while a threshold chosen after inspecting
calibration outcomes requires selection-aware uncertainty accounting.

We introduce \emph{Independently Verified Agent Diagnosis} (\method), a
protocol that makes these trust boundaries explicit. Each diagnosis contains a
bounded causal chain whose claims resolve to immutable trajectory fields and
canonical hashes. A deterministic channel enforces identity, integrity,
temporal, structural, scope, and evidence-coverage constraints; a separately
controlled semantic channel re-evaluates relevance, sufficiency,
counter-evidence, and alternative causal accounts without access to evaluator
labels or the diagnoser's hidden reasoning. Verification can accept, request at
most one auditable revision, abstain, or defer to a human decision.

IVAD also formalizes acceptance at the level where its risk statement is
defined. Reports that share a preregistered task-template, repository, and
environment tuple form one group. The protocol freezes the loss, scoring rule,
finite candidate family, minimum accepted-group count, and tie-breaker before
calibration. Simultaneous one-sided exact-binomial bounds identify candidates
whose group-selective risk is at most a target $\alpha$ under independent-group
and frozen-pipeline assumptions. The selected candidate is evaluated once on
untouched test data; if no candidate is feasible, the protocol returns no
operating point rather than relaxing the target.

\system is the open-source reference implementation of these ideas. It
provides six versioned public contracts, strict trace projection, deterministic
and optional semantic verification, one bounded revision, human authority,
recoverable SQLite state, API and command-line delivery, LangGraph and AutoGen
adapters, two controlled agent labs, and content-addressed evaluation
artifacts. The checked-in validation traverses the complete 24-cell offline
domain--framework--condition matrix with no provider or GPU calls. This
evidence tests executability, parity, provenance, and recovery; it does not
substitute fake-provider behavior for a semantic-effectiveness result.

The v0.7 formal evidence bundle adds a separate, aggregate DeepSeek-only
observation. Its evaluated-results manifest is bound by SHA-256
\nolinkurl{bc09f1b134de9370a3b5209fa5e959bce01abbcdf05c8456af1f069fc4cd3088}.
The run scheduled and evaluated 2,148 plans with no missing cells, completed
B0--B3, and skipped B4/B5 by policy. This report treats the resulting B3
risk--coverage differences as observations within the declared matrix; they do
not identify a causal intervention, establish cross-model behavior, or certify
target risk.

\begin{anonsuppress}
Source code: \href{https://github.com/naturaljam/SpanVouch}
{\nolinkurl{github.com/naturaljam/SpanVouch}}.
\end{anonsuppress}

This Technical Report makes four protocol and engineering contributions:
\begin{itemize}
    \item We define IVAD, a diagnosis protocol whose causal claims bind to
    immutable trajectory fields through a machine-checkable claim--evidence
    contract.
    \item We separate hard factual eligibility from controlled semantic
    verification and make diagnoser--verifier information overlap an explicit
    experimental variable.
    \item We formulate accepted-diagnosis correctness as group-selective risk
    and give a finite-candidate, simultaneous calibration protocol with an
    explicit no-feasible-point outcome.
    \item We develop SpanVouch, a reproducible implementation that connects
    contracts, diagnosis, verification, bounded review, durable recovery,
    framework adapters, and manifest-bound evaluation in one system.
\end{itemize}
